\documentclass{article}
\usepackage[utf8]{inputenc}
\setcounter{secnumdepth}{0} % Set the numbering depth to 0

\title{Autonomous Vehicle Training in CARLA}

\author{Alex Johnson, Julian Rodriguez}

\date{\today}

\begin{document}

\maketitle

\section{Introduction}
Below is a list of the processes used and lessons learned 
as we trained our autonomous vehicle in CARLA. After explaining
our choices and exactly what went right and wrong, we will 
provide general recommendations for autonomous training based
on our experiences

\section{What Went Wrong}

\subsection{Custom PPO vs SB3}
\subsection{Driving Rules}
Our final car appeared to behave erratically at first, as though it wasn't
properly obeying driving rules like staying in the speed limit or using the
right lanes
\paragraph{}

\subsection{?}
\subsection{?}
\subsection{?}

\section{What Went Right}

\subsection{
WandB
}
\paragraph{}{}
Weights and Biases (WandB) is a platform for storing model metrics during
machine learning training. It uses a simple Python API that can be inserted
into a training script. It is free to use and helps us track the model's
losses, collision rate, episode length, and style adherence, create
visuals, and monitor our runtime as we train, without even rendering
the car. It provided a quantitative way to track training, with live
updates and remote access via the website. 
\subsection{Style Vectors}
\paragraph{}
The style vector idea performed very strongly when rendered. Small changes 
to the vectors, which are constrained to the interval [0, 1], resulted in large 
output fluxuations. We found that keeping speed minimal and safety maximal led
to the smoothest car. In the future, training might want to reduce the loss 
related to the style vectors to make the model fluctuate less, ie. make the 
style changes less extreme.
\subsection{?}
\subsection{?}
\subsection{?}

\section{Conclusion}

1 paragraph

\end{document}
